\section{Identification by transport}
\label{sec:identification}

The evolution constructed above gives a family of GGC distributions.
We show that its weak equation determines their Laplace transforms in the same way as deterministic powering.
A logarithmic change of the value variable then identifies the entire family.

\begin{theorem}[Identification of the weak evolution]
\label{thm:identification}
Let $B_0>0$ and $0\le T<\infty$.
Let $(F_t)_{0\le t\le T}$ be a narrowly continuous curve of probability measures on $\R$ satisfying
\begin{equation}
 K_T=\sup_{0\le t\le T}\int_\R y^2F_t(\dd y)<\infty.
 \label{eq:id-assumptions}
\end{equation}
Set $B_t=B_0e^{-t}$ and $U_t=B_t\exp_*F_t$.
Assume that the operator $\mathcal G$ defined in \eqref{eq:generator} satisfies, for every $t\in[0,T]$ and every $\varphi\in C_c^2(\R)$,
\begin{equation}
 F_t(\varphi)-F_0(\varphi)
 =\int_0^t F_u(\mathcal G_{B_u,F_u}\varphi)\,\dd u.
 \label{eq:id-weak-rate}
\end{equation}
Then each $U_t$ is an admissible finite Thorin measure.
Let $\mu_t$ be its zero-drift GGC law.
If $X_0$ has law $\mu_0$, then
\begin{equation}
 \mu_t=\Law(X_0^{e^t}),\qquad 0\le t\le T.
 \label{eq:id-conclusion}
\end{equation}
\end{theorem}

The hypotheses are exactly those supplied by Theorem~\ref{thm:evolution}.
The main analytic point is integration of the resolvent evolution down to the zero Laplace argument.
We establish a uniform logarithmic moment for this purpose, and then obtain the transport equation by approximation of its test functions.

\begin{proof}
For $T=0$, admissibility follows from \eqref{eq:log-moment-admissibility} and the distributional identity is immediate.
We may therefore assume $T>0$.
We distinguish the log-rate probabilities $F_t$ from the value distributions $\mu_t$ throughout the argument.
Lemma~\ref{lem:bounds}, on the compact mass interval $I=[B_0e^{-T},B_0]$, gives a constant $C_T<\infty$ such that
\begin{equation}
 |a_{B_t,F_t}(y)-y|\le C_T,\qquad 0\le k_{B_t,F_t}(y,v)\le1,
 \qquad m_2:=\int_\R v^2\nu_0(\dd v)<\infty.
 \label{eq:id-generator-bounds}
\end{equation}
The compensation in $\mathcal G$ is over the whole jump line.

\medskip\noindent\emph{Admissibility and transform continuity.}
For $s>0$, define
\begin{equation}
 \begin{aligned}
 \ell_s(y)&=\log(1+se^{-y}),&
 \Psi_t(s)&=B_tF_t(\ell_s),\\
 L_t(s)&=\exp\{-\Psi_t(s)\},&
 g_t(s)&=B_t\int_\R\frac{F_t(\dd y)}{s+e^y}.
 \end{aligned}
 \label{eq:id-transforms}
\end{equation}
The elementary bound
\begin{equation}
 0\le\ell_s(y)\le\log(1+s)+y_-
 \label{eq:id-log-test-bound}
\end{equation}
shows that $\int\log(1+1/b)\,U_t(\dd b)<\infty$.
Consequently $U_t$ defines the asserted GGC law $\mu_t$, with transform $L_t$, and differentiation at positive $s$ gives $g_t=\partial_s\Psi_t=-\partial_s\log L_t$.
In particular, no first moment of the positive rate $e^y$ is needed.

The moment assumption implies, uniformly in $t$,
\begin{equation}
 F_t\{|y|>R\}\le\frac{K_T}{R^2},\qquad
 \int_{|y|>R}|y|F_t(\dd y)\le\frac{K_T}{R},\qquad R>0.
 \label{eq:id-rate-tails}
\end{equation}
Truncating the continuous function $\ell_s$, whose growth is at most linear, therefore shows that $t\mapsto\Psi_t(s)$ and $t\mapsto L_t(s)$ are continuous.
For each fixed $t$, dominated convergence, with dominating function $\ell_1$ when $s\le1$, gives $\Psi_t(s)\downarrow0$ as $s\downarrow0$.
Thus $L_t(0)=1$.

\medskip\noindent\emph{Logarithmic moments of the value distributions.}
The following estimate controls both the behavior near zero and the upper tail of the value distributions:
\begin{equation}
 \sup_{0\le t\le T}\int_{(0,\infty)}|\log x|\,\mu_t(\dd x)
 \le C_{\log,T}<\infty,
 \label{eq:id-log-moment}
\end{equation}
where the constant depends only on $B_0,T,K_T$.
This is the integrability needed at the zero Laplace endpoint; an $|x\log x|$ moment is not assumed.

For each fixed $t$, let $X_t$ have distribution $\mu_t$; no joint process $(X_t)$ is assumed.
For the negative part, Lemma~\ref{lem:gamma-dirichlet}, applied in log-rate coordinates, gives
\[
 X_t\overset{d}=G_{B_t}M_t,\qquad
 M_t=\int_\R e^{-y}Q_t(\dd y),\qquad
 Q_t\sim\DP(B_tF_t),
\]
with $G_{B_t}$ independent of $Q_t$.
The logarithmic integrability already verified makes $M_t$ finite almost surely, and it is strictly positive.
In particular, $\mu_t((0,\infty))=1$.
Moreover, $\E Q_t(|y|)=F_t(|y|)<\infty$.
On this almost-sure integrability event, put $m=\int y\,Q_t(\dd y)$ and integrate $e^{-(y-m)}\ge1-(y-m)$ to obtain $e^mM_t\ge1$.
Hence
\[
 (\log M_t)^-\le\left(\int y\,Q_t(\dd y)\right)^+
 \le\int y_+\,Q_t(\dd y),
 \qquad \E(\log M_t)^-\le\sqrt{K_T}.
\]
For the gamma factor,
\[
 \E(\log G_B)^-
 \le\frac1{\Gamma(B)}\int_0^1(-\log x)x^{B-1}\,\dd x
 =\frac1{\Gamma(B)B^2}.
\]
This is bounded uniformly for $B\in I$.
Since $(\log X_t)^-\le(\log G_{B_t})^-+(\log M_t)^-$ under this coupling, the negative part of \eqref{eq:id-log-moment} is controlled.

For the positive part, Tonelli's theorem and the elementary identity
\[
 \log(1+x)=\int_0^\infty\frac{e^{-r}}r(1-e^{-rx})\,\dd r
 \qquad (x\ge0)
\]
give, using $1-e^{-z}\le z$ for $z\ge0$,
\begin{equation}
 \begin{aligned}
 \E\log(1+X_t)
 &=\int_0^\infty\frac{e^{-r}}r(1-L_t(r))\,\dd r
 \le\int_{(0,\infty)}\mathcal{I}(b)\,U_t(\dd b),\\
 \mathcal{I}(b)&:=\int_0^\infty\frac{e^{-r}}r\log(1+r/b)\,\dd r
 \le4\{1+(\log_+(1/b))^2\}.
 \end{aligned}
 \label{eq:id-positive-log}
\end{equation}
Here the logarithmic identity follows, for example, by writing $(1-e^{-rx})/r=\int_0^x e^{-rz}\,\dd z$ and integrating first in $r$.
We next prove the bound on $\mathcal{I}$ for all positive rates.
If $b\ge1$, then $\mathcal{I}(b)\le1/b$, by $\log(1+r/b)\le r/b$.
If $0<b<1$, put $a=\log(1/b)$.
The part $0<r<b$ is at most one.
On $b<r<1$, use $\log(1+r/b)\le\log2+\log(r/b)$; its integral after replacing $e^{-r}$ by one is at most $a\log2+a^2/2$.
On $r>1$, use
\[
 \log(1+r/b)\le\log(1+r)+a\le r+a,\qquad r^{-1}\le1,
\]
which bounds this part by $e^{-1}(1+a)$.
The sum is at most $4(1+a^2)$.
Thus \eqref{eq:id-positive-log} is bounded by $4B_0(1+K_T)$.
Combining this bound with the negative-part estimate, we obtain \eqref{eq:id-log-moment}; one possible choice is
\[
 C_{\log,T}=
 \sup_{B\in I}\frac1{\Gamma(B)B^2}
 +\sqrt{K_T}+4B_0(1+K_T).
\]

The family $(\mu_t)$ is narrowly continuous on $(0,\infty)$.
Indeed, \eqref{eq:id-log-moment} bounds its mass outside $[e^{-R},e^R]$ by $C_{\log,T}/R$.
If $t_n\to t$, tightness and diagonal selection of distribution functions give a weakly convergent subsequence, and this tail bound ensures that its limit is a probability on $(0,\infty)$.
Its transform is $L_t$, by the transform continuity proved above.
These transforms determine a probability uniquely: push the measure forward by $x\mapsto e^{-x}$ to $[0,1]$; the transform values at positive integers give all its positive integer moments, and its zeroth moment is one.
Uniform polynomial approximation of continuous functions then gives uniqueness.
The polynomial approximation argument used below also proves this uniform approximation assertion.
Every subsequential limit is therefore $\mu_t$, proving narrow continuity.

\medskip\noindent\emph{Extension of the log-rate test class.}
The resolvent tests used in the next step are not compactly supported.
We therefore first justify their use in the weak equation.
The weak equation \eqref{eq:id-weak-rate} extends to every $\varphi\in C^2(\R)$ with at most linear growth and with bounded first and second derivatives.
Choose $\chi\in C_c^\infty(\R)$ equal to one on $[-1,1]$ and zero outside $[-2,2]$, and set $\varphi_R(y)=\chi(y/R)\varphi(y)$ for $R\ge1$.
The functions $\varphi_R'$ and $\varphi_R''$ are bounded uniformly in $R$: in the product rule the extra factors $R^{-1}$ and $R^{-2}$ multiply a function of at most linear growth on $|y|\le2R$.
Taylor's formula with full linear compensation therefore gives
\begin{equation}
 \begin{aligned}
 |\mathcal G_{B_t,F_t}\varphi_R(y)|
 &\le (|y|+C_T)|\varphi_R'(y)|
       +\frac{m_2}{2}\|\varphi_R''\|_\infty
 \le C_{T,\varphi}(1+|y|),\\
 \mathcal G_{B_t,F_t}\varphi_R(y)
 &\longrightarrow\mathcal G_{B_t,F_t}\varphi(y).
 \end{aligned}
 \label{eq:id-domain-extension}
\end{equation}
For the second assertion, the compensated jump remainder is bounded by a constant times $v^2$ for every $v$.
The finite second moment in \eqref{eq:id-generator-bounds} therefore controls both the singularity at zero and the tails.
Dominated convergence in $F_t$ and then in time, using $\sup_tF_t(|y|)\le\sqrt{K_T}$, proves the extended weak equation.
The left side converges by the same linear-growth bound.
Also \eqref{eq:id-rate-tails} makes $F_t(\varphi)$ continuous in time.
The extended equation is thus an absolutely continuous scalar identity.

In particular, this applies to
\[
 \varphi_s(y)=(s+e^y)^{-1},\qquad \ell_s(y)=\log(1+se^{-y}),\qquad s>0.
\]
The first function and its first two derivatives are bounded.
For the second, $|\ell_s'|\le1$ and $|\ell_s''|\le1/4$.
Thus the logarithmic test defining $\Psi_t$ is legitimate as well as the resolvent used below.

\medskip\noindent\emph{The resolvent equation.}
Let $X$ denote the coordinate variable on $(0,\infty)$, and write $\E_t H(X)=\int H(x)\mu_t(\dd x)$.
For an integrable test under the exponentially tilted law, set
\[
 \E_{t,s}H(X)=\frac{\int H(x)e^{-sx}\,\mu_t(\dd x)}{L_t(s)}
 \qquad(s>0).
\]
Thus $\E_{t,s}$ is expectation under the normalized exponential tilt of $\mu_t$.
Put
\begin{equation}
 \begin{aligned}
 A_t(s)&=\E_{t,s}[X\log X],\\
 h_t(s)&=\left.\frac{\partial}{\partial q}
 \left[-\partial_s\log\E_t e^{-sX^q}\right]\right|_{q=1}
 =\frac{\dd}{\dd s}\{sA_t(s)\}.
 \end{aligned}
 \label{eq:id-current-tangent}
\end{equation}
Here the power derivative is taken at the current law $\mu_t$, to which Lemma~\ref{lem:tangent} applies.
The last identity can also be obtained by differentiating the normalized exponential tilt.
All derivatives involved are bounded functions of $x$ after multiplication by $e^{-sx^q}$ when $s$ and $q$ range over compact positive intervals.

By Lemma~\ref{lem:generator}, the normalized resolvent identity is
\[
 F_t(\mathcal G_{B_t,F_t}\varphi_s)=\frac{h_t(s)+g_t(s)}{B_t}.
\]
Since $g_t(s)=B_tF_t(\varphi_s)$, the extended weak equation and $B_t'=-B_t$ imply
\begin{equation}
 g_t(s)-g_0(s)=\int_0^t h_u(s)\,\dd u,
 \qquad s>0.
 \label{eq:id-resolvent-evolution}
\end{equation}
The decay of the total Thorin mass thus cancels the extra $g_t(s)$ in the generator identity.

\medskip\noindent\emph{The endpoint at zero.}
To recover the Laplace exponent from its $s$-derivative, we must integrate \eqref{eq:id-resolvent-evolution} down to $s=0$.
Both absolute integrability and the boundary value are needed to recover the normalized Laplace exponent.
Differentiating a normalized tilt in \eqref{eq:id-current-tangent} gives
\begin{equation}
 h_t(s)=A_t(s)-s\E_{t,s}[X^2\log X]
                  +s\E_{t,s}[X]A_t(s).
 \label{eq:id-tangent-expanded}
\end{equation}
We first record joint measurability in time and the Laplace argument.
Fix $0<\sigma<S<\infty$.
If $f_s(x)$ is any one of
\[
 e^{-sx},\quad xe^{-sx},\quad
             x\log x\,e^{-sx},\quad x^2\log x\,e^{-sx},
\]
the map $s\mapsto f_s$ is continuous in the supremum norm on $(0,\infty)$ for $s\in[\sigma,S]$.
Indeed, differentiation in $s$ adds one factor $-x$, and the resulting derivatives are bounded uniformly in $x>0$ and $s\in[\sigma,S]$; the mean value theorem then gives the assertion.
If $(t_n,s_n)\to(t,s)$ in $[0,T]\times[\sigma,S]$, the narrow continuity of the value laws proved above gives
\[
 |\mu_{t_n}(f_{s_n})-\mu_t(f_s)|
 \le\|f_{s_n}-f_s\|_\infty
       +|\mu_{t_n}(f_s)-\mu_t(f_s)|\longrightarrow0.
\]
Since $L_t(s)>0$, division by the jointly continuous denominator $L_t(s)$ preserves continuity.
Equation~\eqref{eq:id-tangent-expanded} therefore shows that $(t,s)\mapsto A_t(s)$ and $(t,s)\mapsto h_t(s)$ are jointly continuous on $[0,T]\times(0,\infty)$, and in particular jointly Borel.
We handle the endpoint $s=0$ by the integral estimates below.

Also, $t\mapsto\mu_t$ is a Borel probability kernel: evaluation on an open set is measurable by bounded continuous approximation, and the monotone-class theorem extends this to every Borel set.
Thus nonnegative Borel integrands may be integrated against $\dd t\,\dd s\,\mu_t(\dd x)$ using Tonelli's theorem.
This supplies the measurability needed for the estimates below, including those with the unbounded test $|\log x|$.

Fix $s_0>0$.
Equations \eqref{eq:id-transforms} and \eqref{eq:id-log-test-bound} give the uniform lower bound
\begin{equation}
 D:=\exp\{-B_0[\log(1+s_0)+\sqrt{K_T}]\}
 \le L_t(s_0)\le L_t(s),\qquad 0<s\le s_0.
 \label{eq:id-transform-lower-bound}
\end{equation}
For every $x>0$,
\[
 \int_0^{s_0}xe^{-sx}\,\dd s\le1,\qquad
 \int_0^{s_0}sx^2e^{-sx}\,\dd s\le1,\qquad
 sx e^{-sx}\le e^{-1}.
\]
Taking absolute values before applying Tonelli therefore yields
\begin{equation}
 \begin{aligned}
 \int_0^{s_0}|A_t(s)|\,\dd s
       &\le D^{-1}\E_t|\log X|,\\
 \int_0^{s_0}s\E_{t,s}[X^2|\log X|]\,\dd s
       &\le D^{-1}\E_t|\log X|,\\
 s\E_{t,s}X&\le(eD)^{-1},\\
 \int_0^T\int_0^{s_0}|h_t(s)|\,\dd s\,\dd t
       &\le T\left(\frac2D+\frac1{eD^2}\right)C_{\log,T}<\infty.
 \end{aligned}
 \label{eq:id-absolute-fubini}
\end{equation}
The last bound includes the product term in \eqref{eq:id-tangent-expanded}; it follows by multiplying the first bound by $(eD)^{-1}$.
Together with the joint measurability just proved, this establishes that $h(t,s):=h_t(s)$ belongs to $L^1([0,T]\times(0,s_0))$.
Fubini's theorem therefore applies on every subrectangle $[0,t]\times(0,s)$ with $t\le T$ and $s\le s_0$.

For every fixed $t$,
\begin{equation}
 \lim_{s\downarrow0}sA_t(s)=0,
 \qquad |sA_t(s)|\le\frac{\E_t|\log X|}{eD}\quad(0<s\le s_0).
 \label{eq:id-zero-endpoint}
\end{equation}
Indeed, $sx e^{-sx}\log x\to0$ pointwise, and its absolute value is at most $|\log x|/e$.
Dominated convergence under $\mu_t$, followed by division by $L_t(s)\ge D$, proves the limit.
The bound in \eqref{eq:id-log-moment} also permits passage of this limit through a time integral.
The pointwise limit in $t$ and this integrable bound suffice for that passage.

The fundamental theorem of calculus applied to \eqref{eq:id-current-tangent} on $[\varepsilon,s]$, followed by \eqref{eq:id-zero-endpoint}, gives
\[
 \int_0^s h_t(r)\,\dd r=sA_t(s).
\]
The integral is absolutely convergent, and \eqref{eq:id-absolute-fubini} justifies its time--Laplace Fubini exchange.
The logarithmic moment in \eqref{eq:id-log-moment} has supplied all the required endpoint control.

\medskip\noindent\emph{The Laplace equation.}
By Tonelli's theorem, $\int_0^s g_t(r)\,\dd r=\Psi_t(s)$.
Integrating \eqref{eq:id-resolvent-evolution}, with the absolute bounds just proved, therefore yields
\begin{equation}
 \Psi_t(s)-\Psi_0(s)=\int_0^t sA_u(s)\,\dd u.
 \label{eq:id-exponent-evolution}
\end{equation}
The normalization is fixed by $\Psi_t(0)=0$ and the vanishing boundary term in \eqref{eq:id-zero-endpoint}.
Applying the absolutely continuous chain rule to $L_t=e^{-\Psi_t}$ gives
\begin{equation}
 L_t(s)-L_0(s)
 =-\int_0^t s\int_{(0,\infty)}x\log x\,e^{-sx}\,\mu_u(\dd x)\,\dd u,
 \qquad s>0.
 \label{eq:id-laplace-evolution}
\end{equation}
Here $L_u(s)A_u(s)=\E_u[X\log X e^{-sX}]$ has removed the tilt normalization.
For each fixed $s>0$, the function $x\log x\,e^{-sx}$ is bounded and continuous and tends to zero at both endpoints.
Narrow continuity of $(\mu_u)$ makes the time integrand in \eqref{eq:id-laplace-evolution} continuous.

\medskip\noindent\emph{The weak equation on value space.}
We next derive the weak equation for the value distributions.
We approximate a test and its derivative simultaneously, so that the transport term converges as well.
Let $H\in C_c^1((0,\infty))$.
On $0<u<1$, set $J(u)=H(-\log u)$.
Because $H$ vanishes near zero and infinity, $J$ extends by zero to a $C^1$ function on $[0,1]$, vanishing near both endpoints.
There are polynomials $p_n$ satisfying
\begin{equation}
 \|p_n-J\|_{\infty,[0,1]}+
 \|p_n'-J'\|_{\infty,[0,1]}\longrightarrow0.
 \label{eq:id-c1-polynomials}
\end{equation}
We construct these polynomials as follows.
For each integer $n\ge1$, form the Bernstein polynomial of $J'$,
\[
 q_n(u)=\sum_{j=0}^n J'(j/n)\binom{n}{j}u^j(1-u)^{n-j}.
\]
For a binomial random variable $N$ with parameters $(n,u)$, their error is bounded by $\E|J'(N/n)-J'(u)|$.
On $|N/n-u|\le\delta$ use the modulus of continuity of $J'$; on its complement use $2\|J'\|_\infty$ and
\[
 \mathbb P\{|N/n-u|>\delta\}
 \le\frac{u(1-u)}{n\delta^2}\le\frac1{4n\delta^2}.
\]
First choosing small $\delta$ and then large $n$ proves uniform convergence of $q_n$ to $J'$.
The polynomials $p_n(u)=J(0)+\int_0^u q_n(v)\,\dd v$ prove \eqref{eq:id-c1-polynomials}.
The same Bernstein estimate applied to any continuous function, without differentiating it, proves the uniform polynomial approximation used above for transform uniqueness.

Put $H_n(x)=p_n(e^{-x})$ and define the value-space transport operator by
\[
 \mathscr{B}H(x)=x\log x\,H'(x),\qquad x>0.
\]
In the following estimates, norms of $H_n,H$ are taken on $(0,\infty)$, whereas norms of $p_n,J$ and their derivatives are taken on $[0,1]$:
\begin{equation}
 \begin{aligned}
 \|H_n-H\|_\infty&\le\|p_n-J\|_\infty,\\
 \|\mathscr{B}H_n-\mathscr{B}H\|_\infty
 &\le\left(\sup_{x>0}|x\log x|e^{-x}\right)
                   \|p_n'-J'\|_\infty\longrightarrow0.
 \end{aligned}
 \label{eq:id-generator-norm}
\end{equation}
The displayed supremum is finite, because its integrand is continuous and vanishes at zero and infinity.
Each $H_n$ is a finite linear combination of a constant and the functions $e^{-jx}$ for positive integers $j$.
Its weak evolution is therefore supplied by \eqref{eq:id-laplace-evolution}; the constant has zero evolution since each $\mu_u$ is a probability.
Both uniform limits in \eqref{eq:id-generator-norm} can be passed through probability integrals and the finite time integral.
We obtain
\begin{equation}
 \mu_t(H)-\mu_0(H)
 =\int_0^t\mu_u(x\log x\,H'(x))\,\dd u,
 \qquad H\in C_c^1((0,\infty)).
 \label{eq:id-value-weak}
\end{equation}
The simultaneous convergence in \eqref{eq:id-generator-norm} is what transfers the Laplace equation to the transport operator.

\medskip\noindent\emph{Transport of logarithmic values.}
The gamma--Dirichlet representation shows that $\mu_t$ gives full mass to $(0,\infty)$.
Define $\lambda_t=(\log)_*\mu_t$.
Thus $\lambda_t$ is the law of the logarithmic value, while $F_t$ continues to describe logarithmic rates.
It is narrowly continuous, since logarithm is a continuous map on $(0,\infty)$, and \eqref{eq:id-log-moment} gives a uniform first moment of $\lambda_t$.

For $\zeta\in C_c^1(\R)$, the function $H(x)=\zeta(\log x)$ lies in $C_c^1((0,\infty))$.
Substitution into \eqref{eq:id-value-weak} gives
\begin{equation}
 \lambda_t(\zeta)-\lambda_0(\zeta)
 =\int_0^t\lambda_u(z\zeta'(z))\,\dd u.
 \label{eq:id-log-transport}
\end{equation}
The characteristics of this equation are dilations.
We verify the resulting identity directly with backward test functions.
Fix $0<t\le T$ and $\zeta\in C_c^\infty(\R)$, and set
\begin{equation}
 \zeta_u(z)=\zeta(e^{t-u}z),\qquad 0\le u\le t.
 \quad\text{Then}\quad
 \partial_u\zeta_u+z\partial_z\zeta_u=0.
 \label{eq:id-backward-test}
\end{equation}
All these tests and their derivatives have support in one compact interval, and their time and spatial derivatives are jointly continuous.

To use these time-dependent tests in the fixed-test identity, we write the increments on a partition.
Take a partition $0=t_0<\cdots<t_N=t$ and decompose each increment as
\[
 \begin{aligned}
 &\lambda_{t_{i+1}}(\zeta_{t_{i+1}})
       -\lambda_{t_i}(\zeta_{t_i})\\
 &\quad=\lambda_{t_{i+1}}(\zeta_{t_{i+1}}-\zeta_{t_i})
       +\lambda_{t_{i+1}}(\zeta_{t_i})
       -\lambda_{t_i}(\zeta_{t_i})\\
 &\quad=\int_{t_i}^{t_{i+1}}
          \lambda_{t_{i+1}}(\partial_u\zeta_u)\,\dd u
       +\int_{t_i}^{t_{i+1}}
          \lambda_u(z\partial_z\zeta_{t_i})\,\dd u,
 \end{aligned}
\]
where the second integral uses the fixed-test identity \eqref{eq:id-log-transport}.
Joint continuity, common compact support, and narrow continuity of $\lambda$ imply joint continuity of $(r,u)\mapsto\lambda_r(\partial_u\zeta_u)$ and of the analogous spatial-derivative pairing.
For example, joint continuity follows by adding a supremum-norm difference of the test functions to the weak continuity term for one fixed test.
On the compact time square it is uniform.
Consequently, as the mesh tends to zero, the summed integrals converge to
\[
 \int_0^t\lambda_u(\partial_u\zeta_u+z\partial_z\zeta_u)\,\dd u=0.
\]
The telescoped left side is $\lambda_t(\zeta)-\lambda_0(\zeta(e^t\,\cdot))$.
Thus
\[
 \lambda_t(\zeta)=\lambda_0(\zeta(e^t\,\cdot)),
 \qquad \zeta\in C_c^\infty(\R).
\]
Compactly supported smooth tests determine finite Borel measures on $\R$, as follows by approximating compactly supported continuous functions uniformly by smooth ones and then approximating interval indicators.
Therefore
\begin{equation}
 \lambda_t=(z\mapsto e^t z)_*\lambda_0,
 \qquad
 \mu_t=(x\mapsto x^{e^t})_*\mu_0=\Law(X_0^{e^t}).
 \label{eq:id-transport-identification}
\end{equation}
The value laws are supported on $(0,\infty)$, so exponentiating the identified logarithmic laws proves \eqref{eq:id-conclusion} for every $t\in[0,T]$.
\end{proof}
